\documentclass[11pt]{article}
\usepackage[a4paper,margin=27mm]{geometry}
\usepackage[T1]{fontenc}
\usepackage[utf8]{inputenc}
\usepackage{lmodern}
\usepackage{amsmath,amssymb,amsthm,mathtools}
\usepackage{booktabs}
\usepackage{array}
\usepackage{xcolor}
\usepackage{microtype}
\usepackage{hyperref}
\hypersetup{colorlinks=true,linkcolor=blue!45!black,citecolor=blue!45!black,urlcolor=blue!55!black}

\newtheorem{theorem}{Theorem}[section]
\newtheorem{proposition}[theorem]{Proposition}
\newtheorem{corollary}[theorem]{Corollary}
\newtheorem{lemma}[theorem]{Lemma}
\newtheorem{remark}[theorem]{Remark}
\newcommand{\Fp}{\mathbf F_p}
\newcommand{\Q}{\mathbf Q}
\newcommand{\Status}[1]{\par\smallskip\noindent\textbf{Status: #1.}\ }
\newcommand{\Res}{\operatorname{Res}}
\newcommand{\Norm}{\operatorname{N}}

\title{\textbf{Arithmetic Quantum Cat Maps:}\\
Exact Fourth-Moment Reductions, Direct Cayley Energy Bounds,\\
and the Remaining BGS Frontier}
\author{Maksim Churilov}
\date{Research draft --- 4 September 2026}

\begin{document}
\maketitle

\begin{abstract}
We develop an exact finite-field reduction of a fourth-moment problem arising
in arithmetic quantum cat maps and isolate the arithmetic obstruction that
remains before any microscopic Bohigas--Giannoni--Schmit (BGS) statement can
be contemplated.  The pair-key geometry gives an exact constrained-pair count
$(K-1)(K-2)$ and a corrected fourth-moment difference distribution
\[
B(t)=O(t)+2(p-2)C_H(t)+(p-2)^2\mathbf 1_{t=0}.
\]
A Cayley transform turns the mixed term into the nonzero additive-difference
energy of a M\"obius image $A=c(H\setminus\{1\})$ of a multiplicative subgroup.
For odd prime subgroup order $K$, Chowla's rational independence theorem for
cotangents gives a characteristic-zero classification of all equal-difference
collisions, and every cleared collision polynomial satisfies
\[
\Res(\Phi_K,P)=K^3M^2,\qquad M\in\mathbf Z.
\]
An explicit slice has
$\Res(\Phi_K,X^2+3X+1)=F_K^2$, so the square factor can be exponentially large.
Nevertheless its rational Cayley model has exact nonzero difference energy
$2N(4N-3)+32$, and a 28-digit finite-field certificate reproduces the same
excess $32$ at $K=131$.

The mixed term can also be controlled directly, without excluding exceptional
primes.  For every subgroup of order
\[
100<K<\frac13p^{3/4},
\]
a source-exact Stepanov estimate on an absolutely irreducible bilinear fiber
gives
\[
E_+(A)<34K^{8/3}.
\]
For odd prime $K$ we further prove a global split-prime budget: if
$\Delta_K(p)$ denotes the non-forced mixed-energy excess, then
\[
\sum_{p\equiv1\,(K)}\Delta_K(p)\log p
\le (K^3-5K^2+6K+1)\bigl((K-1)\log12-3\log K\bigr).
\]
Thus a polynomial-scale pointwise mixed theorem and an $O(K^4)$ global
exceptional-mass theorem are available.  The sharper centered target
$E_+(A)=(K-1)^4/p+O(K^{2+o(1)})$ remains open.  The off/off elliptic-family
estimate remains a theorem draft, and nothing here proves microscopic BGS/GOE
eigenphase statistics.
\end{abstract}

\section{Scope and evidence discipline}

Throughout, $p$ is an odd prime and $H\le \Fp^\times$ is a multiplicative
subgroup of order $K$.  Division is in $\Fp$.  We distinguish:
\begin{itemize}
\item \textbf{Exact}: a current proof with explicit hypotheses;
\item \textbf{Finite certificate}: an exhaustive or kernel-checked statement
      at explicit finite parameters;
\item \textbf{Finite regression}: an implementation check that supports, but
      does not prove, a universal statement;
\item \textbf{Numerical}: descriptive floating-point evidence only;
\item \textbf{Draft/Open}: a named closure condition remains.
\end{itemize}

The distinction between moment information and microscopic spectral statistics
is essential.  Arithmetic cat-map literature contains strong matrix-element,
eigenfunction and quantum-ergodicity results
\cite{KR2005,KRR2007,DJ2024,KORS2025,BS2026}, but these are not interchangeable
with a local eigenphase law.  Conversely, rigorous BGS-type results are known
for special model classes such as mixing quantum graphs under their own
large-graph hypotheses \cite{PW2014,PW2015}; they do not transfer to the
arithmetic family treated here without a new argument.

\section{Exact pair-key geometry}

For $x,y\in\Fp\setminus\{0,1\}$ define
\begin{equation}\label{eq:keydef}
u=\frac{x}{y},\qquad
v=\frac{x-1}{y-1},\qquad d=x-y.
\end{equation}

\begin{proposition}[Off-diagonal pair-key bijection]\label{prop:key}
If $x\ne y$, then $u,v\in\Fp^\times\setminus\{1\}$, $u\ne v$, and
\begin{equation}\label{eq:inverse}
y=\frac{1-v}{u-v},\qquad
x=\frac{u(1-v)}{u-v},\qquad
d=\frac{(u-1)(v-1)}{v-u}.
\end{equation}
Conversely every such ordered pair $(u,v)$ recovers a unique admissible
off-diagonal $(x,y)$.  The diagonal $x=y$ collapses to $(u,v,d)=(1,1,0)$ with
multiplicity $p-2$.
\end{proposition}

\begin{proof}
From $x=uy$ and $x-1=v(y-1)$,
$(u-v)y=1-v$, which gives the inverse formulas and then the formula for $d$.
The diagonal statement is immediate.
\end{proof}

\begin{corollary}[Exact constrained-pair count]\label{cor:k2}
If the two constraints are $u,v\in H$, the number of admissible off-diagonal
ordered pairs is
\[
(K-1)(K-2).
\]
\end{corollary}

\begin{proof}
Choose an ordered pair of distinct elements of $H\setminus\{1\}$ and apply
Proposition~\ref{prop:key}.
\end{proof}

\Status{Exact.} Independent Python/C++ checks and a Lean certificate at
$(p,K)=(53,13)$ guard the implementation, but are not the source of the general
proof.

\section{Fourth-moment coupling and finite Fourier identities}

Give each collapsed key $r=(u_r,v_r,d_r)$ multiplicity $n_r$, equal to one off
the diagonal and $p-2$ at $(1,1,0)$. Define
\begin{equation}\label{eq:Bdef}
B(t)=\sum_{r,s}n_rn_s\,
\mathbf 1_H(u_ru_s)\mathbf 1_H(v_rv_s)\mathbf 1_{d_r+d_s=t}.
\end{equation}
Let $O(t)$ be the off/off contribution and
\[
C_H(t)=\#\left\{(u,v)\in(H\setminus\{1\})^2:
 u\ne v,\ \frac{(u-1)(v-1)}{v-u}=t\right\}.
\]
Then the diagonal bookkeeping is exactly
\begin{equation}\label{eq:Bdecomp}
B(t)=O(t)+2(p-2)C_H(t)+(p-2)^2\mathbf 1_{t=0}.
\end{equation}
The coefficient $2$ is load-bearing: diagonal/off and off/diagonal are distinct
ordered classes.

\begin{proposition}[Exact total coupled count]\label{prop:v4}
\[
|V_4|=(K^2+1)p^2-(6K^2+3K)p+13K^2-3K.
\]
\end{proposition}

Let $S=H^\perp$ be the multiplicative-character annihilator, so
$m=|S|=(p-1)/K$. In the maintained twisted-Jacobi convention,
\begin{equation}\label{eq:fourier}
\sum_{\chi,\psi\in S}|J(\chi,\psi;q)|^4=m^2\Sigma_4(q),
\qquad
\Sigma_4(q)=\sum_{t\in\Fp}B(t)e_p(qt).
\end{equation}
Writing $V=\sum_tB(t)$ and $Z=B(0)$,
\begin{equation}\label{eq:meanvar}
\mathbb E_{q\ne0}\Sigma_4(q)=\frac{pZ-V}{p-1},
\qquad
\operatorname{Var}_{q\ne0}\Sigma_4(q)
=p\operatorname{Var}_{t\ne0}B(t).
\end{equation}
These follow directly from finite Fourier orthogonality and Parseval; no
independence heuristic enters.

\section{Cayley transform and additive energy}

Put
\[
c(h)=\frac{h+1}{h-1},\qquad h\in H\setminus\{1\}.
\]
For $u\ne v$,
\begin{equation}\label{eq:cayley}
\frac{(u-1)(v-1)}{v-u}=\frac{2}{c(u)-c(v)}.
\end{equation}
Thus for $A=c(H\setminus\{1\})$ the nonzero multiplicities of $C_H$ are the
nonzero difference multiplicities of $A$, up to the permutation $x\mapsto2/x$.
Consequently
\begin{equation}\label{eq:energy}
\sum_t C_H(t)^2=E_+(A)-|A|^2,
\qquad E_+(A)=\sum_x r_{A-A}(x)^2.
\end{equation}
Since $c(h^{-1})=-c(h)$, $A=-A$ and $0\notin A$. With $n=|A|=K-1$, identity
and inversion-swap alone force the nonzero-energy baseline
\begin{equation}\label{eq:baseline}
n(2n-3)=(K-1)(2K-5).
\end{equation}

A separate exact matching audit of all $6!=720$ integer systems obtained by
clearing and matching Cayley collisions gives ranks
\[
30\ (\text{rank }2),\qquad240\ (\text{rank }3),\qquad450\ (\text{rank }4),
\]
and an extra rank drop can occur only in characteristic $2$. Complete fixed-$K$
cyclotomic enumeration currently gives
\begin{equation}\label{eq:fixed-exceptions}
E_{11}=\{23,67,89\},\qquad
E_{13}=\{53,79,131,157,313\}.
\end{equation}
These are exact computer-assisted finite statements, not growing-$K$ theorems.

\section{Prime-order characteristic-zero rigidity}

Let $K$ now be an odd prime and $\zeta=e^{2\pi i/K}$. For $1\le j\le K-1$,
\begin{equation}\label{eq:cot}
c(\zeta^j)=-i\cot\frac{\pi j}{K}.
\end{equation}
Chowla proved that
\[
\cot\frac{\pi a}{K},\qquad1\le a\le\frac{K-1}{2},
\]
are linearly independent over $\Q$ \cite{Chowla1970}; Wang states this prime
case explicitly and extends it \cite{Wang1982}.

\begin{lemma}[Cross-polytope equal differences]\label{lem:crosspoly}
Let $x,y,u,v$ be vertices of $\{\pm e_1,\ldots,\pm e_m\}$ with $x\ne y$ and
$u\ne v$. If $x-y=u-v$, then either $(u,v)=(x,y)$ or
$(u,v)=(-y,-x)$.
\end{lemma}

\begin{proof}
The equality is $x+v=u+y$. If the common sum vanishes, then
$v=-x$ and $u=-y$. Otherwise a nonzero sum of two signed basis vectors
uniquely determines its unordered pair of summands. The alternative pairing
would force $x=y$, excluded.
\end{proof}

\begin{theorem}[Prime-order collision rigidity]\label{thm:rigidity}
Every characteristic-zero equal-difference collision in
$\{c(\zeta^j):1\le j<K\}$ is either identity or central inversion-swap.
\end{theorem}

\begin{proof}
Use \eqref{eq:cot}, replace each cotangent with the corresponding signed basis
vector via Chowla's rational independence, and apply
Lemma~\ref{lem:crosspoly}.
\end{proof}

\section{Square cyclotomic resultants}

For nonzero residues $i,j,r,s\pmod K$, with $i\ne j$ and $r\ne s$, choose
integer lifts and define the cleared collision polynomial
\begin{align}\label{eq:Pdef}
P(X)={}&X^{r+s+i}-X^{r+s+j}-X^{r+i+j}+2X^{r+j}-X^r\\
&+X^{s+i+j}-2X^{s+i}+X^s+X^i-X^j.\nonumber
\end{align}
Direct denominator clearing gives
\begin{equation}\label{eq:clear}
2P(X)=\bigl(c(X^r)-c(X^s)-c(X^i)+c(X^j)\bigr)
\prod_{a\in\{r,s,i,j\}}(X^a-1).
\end{equation}
The coefficient $\ell^1$ norm satisfies $\|P\|_1\le12$.

\begin{lemma}[Cubic zero and reciprocity]\label{lem:reciprocity}
Writing $T=i+j+r+s$, one has
\[
P(X)=(X-1)^3Q(X),\qquad Q\in\mathbf Z[X],
\]
and
\[
Q(X^{-1})=X^{3-T}Q(X).
\]
\end{lemma}

\begin{proof}
In \eqref{eq:clear}, cancel the denominator of each Cayley term separately;
every summand retains three factors $X^a-1$, proving the cubic zero.
Furthermore $c(X^{-a})=-c(X^a)$ and
$X^{-a}-1=-(X^a-1)/X^a$. The four denominator factors contribute positive
sign and total power $X^{-T}$, while the Cayley-difference contributes one
minus sign. Hence
$P(X^{-1})=-X^{-T}P(X)$, which gives the formula for $Q$.
\end{proof}

\begin{theorem}[Square-resultant theorem]\label{thm:square-resultant}
For every odd prime $K$ and every admissible collision quadruple,
\begin{equation}\label{eq:square-resultant}
\Res(\Phi_K,P)=K^3M^2
\end{equation}
for some $M\in\mathbf Z$. For a non-forced characteristic-zero collision,
$M\ne0$.
\end{theorem}

\begin{proof}
Let $L=\Q(\zeta)$ and $L^+$ be its maximal real subfield. Since $2$ is
invertible modulo $K$, choose $m$ with
$2m\equiv3-T\pmod K$ and put
$\beta=\zeta^mQ(\zeta)$. By Lemma~\ref{lem:reciprocity},
\[
\overline\beta=\zeta^{-m}Q(\zeta^{-1})
=\zeta^{-m+3-T}Q(\zeta)=\beta.
\]
Thus $\beta$ is an algebraic integer in $L^+$. Multiplication by $\zeta^m$
does not change the norm, so
\[
\Norm_{L/\Q}(Q(\zeta))
=\Norm_{L^+/\Q}(\beta)^2=M^2.
\]
Also $\Norm_{L/\Q}(\zeta-1)=\Phi_K(1)=K$. Taking norms of
$P(\zeta)=(\zeta-1)^3Q(\zeta)$ proves \eqref{eq:square-resultant}.
The non-forced case has $P(\zeta)\ne0$ by
Theorem~\ref{thm:rigidity}.
\end{proof}

\begin{corollary}[Improved no-collision cutoff]\label{cor:cutoff}
Suppose $K\mid p-1$. If a non-forced finite-field collision exists, then
\begin{equation}\label{eq:cutoff}
K^3p^2\le 12^{K-1}.
\end{equation}
Consequently $K^3p^2>12^{K-1}$ implies the exact baseline
\eqref{eq:baseline}.
\end{corollary}

\begin{proof}
A finite-field collision makes $P$ and $\Phi_K$ share a root modulo $p$, so
$p\mid\Res(\Phi_K,P)$. Since $K\mid p-1$, $p\ne K$. Equation
\eqref{eq:square-resultant} gives $p\mid M$, hence
$K^3p^2\le|\Res(\Phi_K,P)|$. On the other hand
$|P(\zeta^a)|\le\|P\|_1\le12$ for all $a$, giving the upper bound.
\end{proof}

The cutoff remains exponential in $K$. The next section shows that this is not
merely an artifact of a crude norm estimate.

\section{An explicit Fibonacci collision slice}

Consider
\begin{equation}\label{eq:fibcollision}
c(h^2)-c(h)=c(h)-c(h^{-1}).
\end{equation}
Since $c(h^{-1})=-c(h)$,
\[
c(h^2)-3c(h)
=-\frac{2(h^2+3h+1)}{(h-1)(h+1)}.
\]
For odd-order $h\ne1$ we also have $h\ne-1$, so
\eqref{eq:fibcollision} is equivalent to
\begin{equation}\label{eq:q}
q(h)=h^2+3h+1=0.
\end{equation}

\begin{theorem}[Fibonacci resultant identity]\label{thm:fibres}
For every odd prime $K$,
\begin{equation}\label{eq:fibres}
\Res(\Phi_K(X),X^2+3X+1)=F_K^2.
\end{equation}
For the corresponding cleared Cayley collision polynomial, the full resultant
is $K^3F_K^2$.
\end{theorem}

\begin{proof}
Let $\alpha,\beta$ be the roots of $X^2+3X+1$, so
$\alpha\beta=1$, $\alpha+\beta=-3$, and
$\Phi_K(X)=(X^K-1)/(X-1)$. Therefore
\[
\Res(\Phi_K,q)
=\frac{(\alpha^K-1)(\beta^K-1)}{(\alpha-1)(\beta-1)}
=\frac{2-(\alpha^K+\beta^K)}5.
\]
With $\varphi=(1+\sqrt5)/2$, the roots are $-\varphi^2$ and
$-\varphi^{-2}$. Since $K$ is odd,
$\alpha^K+\beta^K=-L_{2K}$, and the standard identity
$L_{2K}=5F_K^2-2$ for odd $K$ gives \eqref{eq:fibres}. Comparing with
Theorem~\ref{thm:square-resultant}, or directly taking the cleared denominator
norm, gives $K^3F_K^2$.
\end{proof}

Because $F_K\sim\varphi^K/\sqrt5$, Theorem~\ref{thm:fibres} proves that no
bound $|M|\le K^C$ can hold for all collision resultants. Thus any useful
polynomial-scale theorem must exploit eligible prime divisors or directly bound
energy rather than total resultant size.

If $p$ is prime and $K\mid p-1$, then \eqref{eq:fibcollision} occurs over
$\Fp$ if and only if $p\mid F_K$. Indeed $\Phi_K$ splits completely over
$\Fp$, and \eqref{eq:fibres} is exactly the common-root criterion.

\section{Fibonacci/Lucas ratios form an almost-Sidon set}

Assume $q(h)=0$ and put $t=c(h)$. Substitution
$h=(t+1)/(t-1)$ into $q(h)=0$ gives
\begin{equation}\label{eq:tquad}
5t^2=1.
\end{equation}
A direct Binet calculation, equivalently an induction using the standard
Fibonacci/Lucas recurrences, gives
\begin{equation}\label{eq:ratio}
F_n c(h^n)=L_n c(h),\qquad n\ge1.
\end{equation}
Thus the characteristic-zero model for half of the centrally symmetric Cayley
set is
\[
x_n=\frac{L_n}{F_n},\qquad
S_N=\{\pm x_1,\ldots,\pm x_N\}.
\]
Two identities used repeatedly are
\begin{equation}\label{eq:sumdiff}
x_m+x_n=\frac{2F_{m+n}}{F_mF_n},
\qquad
|x_m-x_n|=\frac{2F_{|m-n|}}{F_mF_n}.
\end{equation}

Carmichael's primitive-divisor theorem implies that $F_j$ has a prime divisor
not dividing any earlier Fibonacci number for every $j>12$
\cite{Carmichael1913,Yabuta2001}. We use this only as a divisibility separation
principle; the finite range through $12$ is checked exactly.

\begin{theorem}[Exact almost-Sidon energy]\label{thm:almost-sidon}
For every $N\ge3$, if
\[
E^*(S_N)=\sum_{d\ne0}r_{S_N-S_N}(d)^2,
\]
then
\begin{equation}\label{eq:almost-sidon}
E^*(S_N)=2N(4N-3)+32.
\end{equation}
The excess over the centrally symmetric forced baseline is exactly $32$,
independent of $N$.
\end{theorem}

\begin{proof}
First classify relations between positive $x_n$.

For equal absolute differences, suppose $m>n$, $r>s$ and
\[
\frac{F_{m-n}}{F_mF_n}=\frac{F_{r-s}}{F_rF_s}.
\]
If, say, $m>r$ and $m>12$, a primitive divisor of $F_m$ divides the right
cross-multiplied side but none of $F_{m-n},F_r,F_s$, a contradiction. Thus
beyond the bounded range the upper indices agree. Then
$F_{m-n}/F_n$ is strictly decreasing in $n$, so the lower indices agree.

For equal sums, if $m+n>r+s$ and $m+n>12$, a primitive divisor of
$F_{m+n}$ cannot divide any of $F_{r+s},F_m,F_n$, again a contradiction.
Thus the index sums agree. Assume $m\le n$, $r\le s$, and put
$T=m+n=r+s$. Equality reduces to $F_mF_n=F_rF_s$. The exact product identity
\[
5F_mF_n=L_T-(-1)^mL_{n-m}
\]
then forces the unordered pairs $\{m,n\}$ and $\{r,s\}$ to agree.

Finally $1\le x_n\le3$. A positive sum can equal an absolute difference only
at the endpoint value $2$, giving $x_1+x_1=x_2-x_1$. The remaining bounded
relations are exactly those involving $x_1=1,x_2=3,x_3=2$:
\[
|x_3-x_1|=|x_2-x_3|=1,
\quad x_1+x_1=x_2-x_1=2,
\quad x_1+x_2=x_3+x_3=4.
\]
The finite residual statement is independently kernel-checked in Lean for
indices through $12$.

For a centrally symmetric set of size $2N$, forced inversion pairing alone
contributes $2N(4N-3)$. The two oriented bins $\pm1$ have multiplicity $4$
rather than two independent multiplicity-$2$ bins, contributing excess $8$
each. The four bins $\pm2,\pm4$ have multiplicity $3$ rather than a
multiplicity-$1$ and a multiplicity-$2$ bin, contributing excess $4$ each.
The total excess is $2\cdot8+4\cdot4=32$.
\end{proof}

\begin{remark}
The broader theorem of Bilu--Hanrot--Voutier \cite{BHV2001} supplies primitive
divisors for much more general Lucas/Lehmer sequences. For
Theorem~\ref{thm:almost-sidon}, the classical Fibonacci theorem already
suffices.
\end{remark}

\section{A 28-digit exact finite-field certificate}

Set
\[
K=131,\qquad
p=F_{131}=1066340417491710595814572169.
\]
The repository proves the primality of $p$ by the Lucas criterion using the
complete factorization
\[
p-1=2^3\cdot11\cdot89\cdot131\cdot199\cdot521\cdot2081\cdot9901
\cdot19801\cdot24571
\]
and base $3$: it checks
\[
3^{p-1}\equiv1\pmod p,
\qquad
\gcd\bigl(3^{(p-1)/q}-1,p\bigr)=1
\]
for every prime $q$ in this factorization.

The explicit integer
\[
h=188706367382658195937453923
\]
satisfies
\[
h^2+3h+1\equiv0\pmod p,
\qquad h^{131}\equiv1\pmod p,
\qquad h\ne1.
\]
Since $131$ is prime, $h$ has exact order $131$. Exact arbitrary-precision
computation gives
\begin{equation}\label{eq:giant}
\sum_t C_H(t)^2=33442,
\qquad
(K-1)(2K-5)=33410,
\qquad
\text{excess}=32.
\end{equation}
The maximum nonzero difference multiplicity is $4$, with histogram
\[
1^{\times126},\qquad2^{\times8312},\qquad3^{\times4},\qquad4^{\times2}.
\]
Independent Python and C++20 Boost.Multiprecision implementations verify these
numbers and the primality certificate. The finite low-index relation
classification in Theorem~\ref{thm:almost-sidon} is separately checked by
Lean 4.

\Status{Finite certificate.} This example is not extrapolated to an infinite
sequence of Fibonacci primes. Its role is to show that very large exceptional
split primes can coexist with only constant excess energy.

For comparison, the independent design $(p,K)=(461,23)$ has
\[
\sum_tC_H(t)^2=1190,\qquad\text{baseline}=902,\qquad\text{excess}=288.
\]
Thus exceptional reductions can have different collision geometry.

\section{Direct polynomial-scale Cayley energy}

The Fibonacci obstruction changes the useful closure condition from
``eliminate every extra collision'' to ``control total collision energy''.  In
fact a uniform polynomial-scale theorem follows directly from the sum fibers of
$A$.

For $t\in\Fp$, put
\[
r(t)=\#\{(a,b)\in A^2:a+b=t\}.
\]
If $a=c(x)$ and $b=c(y)$ with $x,y\in H\setminus\{1\}$, clearing the two
Cayley denominators gives
\begin{equation}\label{eq:bilinear-fiber}
P_t(x,y)=(2-t)xy+t(x+y)-(2+t)=0.
\end{equation}

\begin{lemma}[Absolute irreducibility of every sum fiber]\label{lem:bilinear-irred}
For every $t\in\Fp$, the polynomial $P_t$ is absolutely irreducible.
\end{lemma}

\begin{proof}
A nonconstant bilinear polynomial $axy+bx+cy+d$ factors over the algebraic
closure only if $ad-bc=0$.  For \eqref{eq:bilinear-fiber},
\[
(2-t)(-(2+t))-t^2=-4\ne0
\]
because $p$ is odd.  The linear special case $t=2$ is also irreducible.
\end{proof}

\begin{theorem}[Direct Cayley energy bound]\label{thm:direct-energy}
Let $p$ be odd and let $H\le\Fp^\times$ have order $K$.  If
\begin{equation}\label{eq:direct-range}
100<K<\frac13p^{3/4},
\end{equation}
then, for $A=c(H\setminus\{1\})$,
\begin{equation}\label{eq:direct-energy}
E_+(A)<34K^{8/3}.
\end{equation}
Consequently
\begin{equation}\label{eq:mixed-energy-direct}
\sum_t C_H(t)^2<34K^{8/3}.
\end{equation}
\end{theorem}

\begin{proof}
For $t\ne0,-2$, the polynomial \eqref{eq:bilinear-fiber} has bidegree
$(1,1)$,
\[
P_t(0,0)=-(2+t)\ne0,
\qquad
\deg P_t(x,0)=1.
\]
Theorem~2 of Vyugin--Makarychev \cite{VM2015} applies to the subgroup product
$H\times H$ under \eqref{eq:direct-range} and gives
\[
r(t)\le32K^{2/3}.
\]
The theorem actually counts the enlarged fiber in $H^2$, including the
extraneous cleared-denominator point $(1,1)$, so it is a valid upper bound for
$r(t)$.

Since $A=-A$ and $0\notin A$, $r(0)=K-1$.  Trivially $r(-2)\le K-1$.
Using $\sum_t r(t)=|A|^2=(K-1)^2$,
\[
\begin{aligned}
E_+(A)&=\sum_t r(t)^2\\
&\le2(K-1)^2+32K^{2/3}\sum_{t\ne0,-2}r(t)\\
&<34K^{8/3}.
\end{aligned}
\]
Equation \eqref{eq:mixed-energy-direct} follows from \eqref{eq:energy}.
\end{proof}

The load-bearing source is the stated Theorem~2 of the 2015 arXiv version
\cite{VM2015}; the later journal publication is \cite{MV2019}.  In particular,
the source condition is $\deg P(x,0)\ge1$, not a strict inequality.  This was
checked against the primary PDF because dropping the equality bar would make
the present degree-one application invalid.

Let $M(t)=2(p-2)C_H(t)$ denote the mixed summand in \eqref{eq:Bdecomp}.  Since
variance is bounded by the corresponding second moment,
\begin{equation}\label{eq:mixed-var}
\operatorname{Var}_{t\ne0}M(t)
\le\frac{4(p-2)^2}{p-1}\sum_tC_H(t)^2
<136pK^{8/3}.
\end{equation}
Thus a genuine pointwise polynomial-scale mixed bound is proved.  The exponent
$8/3$, however, still loses $K^{2/3}$ relative to the natural quadratic
collision scale.

\section{Centered energy and the exact Fourier/character target}

A raw global estimate $E_+(A)=O(K^2)$ cannot be the correct statement in every
density regime.  With the unnormalized additive Fourier transform
\[
\widehat{1_A}(q)=\sum_{a\in A}e_p(qa),
\]
Parseval gives the exact identity
\begin{equation}\label{eq:centered-energy}
E_+(A)=\frac{|A|^4}{p}
+\frac1p\sum_{q\ne0}|\widehat{1_A}(q)|^4.
\end{equation}
The zero-frequency density term is unavoidable.  The sharp pointwise target is
therefore
\begin{equation}\label{eq:centered-target}
E_+(A)=\frac{(K-1)^4}{p}+O(K^{2+o(1)}),
\end{equation}
or equivalently
\[
\sum_{q\ne0}|\widehat{1_A}(q)|^4\ll pK^{2+o(1)}.
\]
In the intended regime $K=o(p^{1/4})$, the density term is negligible at the
quadratic collision scale.

There is an exact character decomposition.  Put $S=H^\perp$ and
$m=|S|=(p-1)/K$.  With $y=(h-1)^{-1}$,
\[
h=\frac{y+1}{y},\qquad c(h)=1+2y,
\]
and character orthogonality gives
\begin{equation}\label{eq:Tchi}
\widehat{1_A}(q)=\frac{e_p(q)}m\sum_{\chi\in S}T_\chi(q),
\qquad
T_\chi(q)=\sum_{y\ne0,-1}\chi\!\left(\frac{y+1}{y}\right)e_p(2qy).
\end{equation}
The Gram matrix is explicit:
\begin{equation}\label{eq:gram}
\sum_qT_\chi(q)\overline{T_\eta(q)}
=\begin{cases}
p(p-2),&\chi=\eta,\\
-p,&\chi\ne\eta.
\end{cases}
\end{equation}
Indeed additive orthogonality forces the same $y$, while
$y\mapsto(y+1)/y$ bijects $\Fp\setminus\{0,-1\}$ with
$\Fp^\times\setminus\{1\}$.  Equation \eqref{eq:gram} recovers the exact
second moment
\[
\sum_q|\widehat{1_A}(q)|^2=p(K-1),
\]
but the target \eqref{eq:centered-target} is an $L^4$ statement and therefore
requires cancellation beyond this Gram identity.

Equivalently, solving $c(h)-c(k)=t$ gives
\begin{equation}\label{eq:unipotent}
k=u_t(h)=\frac{(2-t)h+t}{2+t-th}.
\end{equation}
The matrix
\[
U_t=\begin{pmatrix}2-t&t\\-t&2+t\end{pmatrix}
\]
has determinant $4$ and trace $4$.  Therefore $\frac12U_t$ lies in
$\mathrm{SL}_2(\Fp)$, has trace $2$, and is unipotent.  In fact $u_t$ is the
Cayley conjugate of the translation $z\mapsto z-t$.  Thus the same problem is
an $L^2$ intersection problem between a multiplicative torus and a
one-parameter unipotent family.

A second exact reformulation relates one Cayley sum fiber to one distinguished
product class of a shifted subgroup.  For $t\ne2$ and
$\lambda=t/(2-t)$,
\[
((2-t)x+t)((2-t)y+t)=(2-t)P_t(x,y)+4,
\]
so on $P_t=0$,
\[
(x+\lambda)(y+\lambda)=(1+\lambda)^2.
\]
Known near-quadratic total multiplicative-energy estimates for shifted
subgroups \cite{MSS2018}, and generic M\"obius incidence bounds such as
\cite{RW2022}, are relevant but do not by themselves control the square-sum of
this moving distinguished product class.  This is the present pointwise gap.

\Status{Open at the centered quadratic scale.}  The pointwise
$K^{8/3}$ theorem is exact; no theorem checked in the current audit upgrades it
to \eqref{eq:centered-target}.

\section{A global split-prime budget for non-forced collision mass}

The resultant structure gives a different kind of improvement when $K$ is an
odd prime: an exact average/sparsity theorem over all split rational primes.
For $p\equiv1\pmod K$, define
\[
\Delta_K(p)=\sum_{d\ne0}r_{A-A}(d)^2-(K-1)(2K-5).
\]
This is exactly the number of collision quadruples not belonging to identity
or central inversion.

Simultaneous multiplication of all four nonzero exponent indices by
$a\in\mathbf F_K^\times$ acts freely on collision quadruples.  The number of
non-forced scaling orbits is
\begin{equation}\label{eq:NK}
N_K=(K-1)(K-2)^2-(2K-5)=K^3-5K^2+6K+1.
\end{equation}
Choose the unique representative of each orbit with first index $i=1$, and
write its collision polynomial and resultant as $P_q$ and $R_q$.

\begin{theorem}[Weighted split-prime collision budget]\label{thm:prime-budget}
For every odd prime $K$,
\begin{equation}\label{eq:prime-budget}
\sum_{p\equiv1\,(K)}\Delta_K(p)\log p
\le N_K\bigl((K-1)\log12-3\log K\bigr),
\end{equation}
where the sum is over rational primes.
\end{theorem}

\begin{proof}
Fix a split prime $p$ and a generator $h$ of its order-$K$ subgroup.  For one
normalized orbit $q$, its scaled member $a q$ is a collision precisely when
$P_q(h^a)=0$.  Thus the number $s_q(p)$ of colliding members of the orbit is
the number of common roots of $P_q$ and $\Phi_K$ among the $K-1$ primitive
$K$-th roots in $\Fp$.

Because $p\ne K$, $\Phi_K$ is separable modulo $p$ and splits completely.
Hensel-lift its simple roots to $\mathbf Z_p$.  Since $\Phi_K$ is monic,
\[
R_q=\prod_{\Phi_K(\alpha)=0}P_q(\alpha),
\]
so every common root contributes at least one factor of $p$ and
\[
s_q(p)\le v_p(R_q).
\]
On the other hand Theorem~\ref{thm:square-resultant} gives
$R_q=K^3M_q^2$.  A split prime is never $K$, hence
$v_p(R_q)=2v_p(M_q)$.  Summing first over split primes and then over normalized
orbits,
\[
\begin{aligned}
\sum_{p\equiv1\,(K)}\Delta_K(p)\log p
&=\sum_q\sum_{p\equiv1\,(K)}s_q(p)\log p\\
&\le\sum_q2\log|M_q|\\
&=\sum_q\log\frac{|R_q|}{K^3}\\
&\le N_K\bigl((K-1)\log12-3\log K\bigr),
\end{aligned}
\]
using $|R_q|\le12^{K-1}$.
\end{proof}

For every $P\ge2$ and $L>0$, Theorem~\ref{thm:prime-budget} immediately gives
\[
\sum_{\substack{p\equiv1\,(K)\\p\ge P}}\Delta_K(p)
\le\frac{N_KC_K}{\log P},
\qquad
\#\{p\ge P:p\equiv1\ (K),\ \Delta_K(p)\ge L\}
\le\frac{N_KC_K}{L\log P},
\]
where $C_K=(K-1)\log12-3\log K$.  No prime number theorem in arithmetic
progressions is used.  A normalized average over a prescribed interval would
need an independent lower bound for the number of split primes in that
varying-modulus interval.

\Status{Exact average/sparsity theorem.}  Exact normalized resultant
regressions for the complete $K=11,13$ exceptional sets verify the
orbit/common-root/$p$-adic-valuation mechanism.  They are finite checks, not
the proof of Theorem~\ref{thm:prime-budget}.

\section{The off/off elliptic route}

After separating the mixed term, the unresolved contribution is the centered
energy of $O(t)$. Historical symbolic work organizes the relevant fibers as a
one-parameter genus-one family and suggests a quartic/Weierstrass model, a
nonconstant rational $j$-map, visible bad fibers and same-$j$ parameter
symmetries.

\Status{Theorem draft.} A publication-grade proof must re-establish or certify
all of the following for the actual family:
\begin{enumerate}
\item a smooth projective model on an explicit open base;
\item the complete discriminant/bad-fiber locus, including characteristic
      exceptions;
\item affine/projective boundary corrections;
\item purity, lissity, geometric irreducibility, and uniform
      conductor/ramification control of the relevant $R^1$ sheaf;
\item all geometric isomorphisms and twists between parameter fibers;
\item separable and inseparable isogeny exclusions or explicit exceptional
      strata;
\item the exact hypotheses of the trace-sheaf correlation theorem used.
\end{enumerate}

Michel-type family arguments \cite{Michel1995}, quantitative sheaf theory
\cite{SFFK2023}, Tate/isogeny technology \cite{Zarhin2013}, and the modern
joint-family example of Grove--Saad \cite{GroveSaad2025} are relevant tools,
but none is a plug-in proof for this pencil.

Conditional on the required square-root correlation with uniform constants,
the existing combinatorics suggest
\begin{equation}\label{eq:offoff-target}
O\!\left(K^2p+K^4p^{1/2}+K^4\right),
\end{equation}
with intended regime $K=o(p^{1/4})$. Equation
\eqref{eq:offoff-target} remains a draft conclusion.

\section{Finite spectral diagnostics and the microscopic BGS gap}

The maintained finite spectral code diagonalizes the full complex Hermitian
sector matrix
\[
\mathcal H(q)=M(q)+M(q)^*.
\]
A reconstruction-era implementation briefly replaced it by its entrywise real
part; that bug was removed and guarded by regression. Current spacing records
are descriptive empirical-CDF comparisons only. Pooled unfolded gaps are
dependent, so no IID Kolmogorov--Smirnov significance is attached. The
forward-window statistic retained in the repository is likewise not stationary
number variance.

The logical chain is
\[
\boxed{\text{exact pair-key/Fourier reduction}}
\longrightarrow
\boxed{\text{mixed: }K^{8/3}\text{ pointwise; centered }K^{2+o(1)}\text{ open}}
\longrightarrow
\boxed{\text{off/off elliptic correlation: draft}}
\longrightarrow
\boxed{\text{moment theorem}}
\longrightarrow
\boxed{\text{microscopic BGS bridge: open}}.
\]
The prime-budget theorem supplies an additional averaged control on exceptional
mixed reductions; it does not silently replace the pointwise box.  Closing an
earlier box does not close a later one.  A local-statistics theorem still
requires an explicit comparison/local-law mechanism, or a genuinely arithmetic
replacement, at mean-spacing scale with the correct symmetry class and
uniformity.

\section{Reproducibility and independent failure modes}

The maintained verification stack is deliberately heterogeneous:
\begin{itemize}
\item Python/SymPy: transparent symbolic algebra, exact resultants,
      Galois-orbit/$p$-adic budget regressions, finite-field reference
      computations, primality and energy certificates;
\item C++20: independent exact/exhaustive kernels, the direct Cayley-energy
      verifier, and the 28-digit Boost.Multiprecision Fibonacci witness;
\item Lean 4: kernel-checked pair-key certificates, fixed Cayley controls,
      torus facts, historical falsifiers, and the bounded Fibonacci-ratio
      residual;
\item TeX/CI: publication-facing integration and independent manuscript build.
\end{itemize}

The C++ arbitrary-precision verifier intentionally uses a different code path
from Python. The Lean module does not formalize Carmichael's theorem; it
certifies exactly the finite range left after the external primitive-divisor
theorem is applied. This separation keeps the evidence graph honest.

The current tree contains no Julia sources. Historical Julia files remain
accounted for in repository provenance; false algorithms were not ported merely
to increase language diversity.

\section{Conclusion}

The exact pair-key/Fourier reduction and Cayley transform isolate the
fourth-moment arithmetic obstruction sharply. Prime-order cotangent rigidity
removes characteristic-zero accidental collisions. The identity
\[
\Res(\Phi_K,P)=K^3M^2
\]
explains the systematic square structure seen in fixed-$K$ computations. The
Fibonacci slice shows why eliminating every exceptional prime is the wrong
target: large collision norms and very large exceptional split primes need not
produce large energy.

The direct bilinear-fiber theorem now closes a genuine pointwise
polynomial-scale part of the mixed problem:
\[
E_+(c(H\setminus\{1\}))<34K^{8/3}
\quad\text{for}\quad
100<K<p^{3/4}/3.
\]
The Galois-orbit resultant argument independently gives a global $O(K^4)$
logarithmic budget for all non-forced collision mass across split primes. What
remains in the mixed term is a sharper centered-energy problem, not the earlier
zero-collision problem. The off/off elliptic correlation remains a separate
geometric theorem draft. Even success at the full moment level would still
leave microscopic BGS as an independent local-statistics problem.

\begin{thebibliography}{99}

\bibitem{KR2005}
P.~Kurlberg and Z.~Rudnick,
\emph{On the distribution of matrix elements for the quantum cat map},
Ann. of Math. \textbf{161} (2005), 489--507,
doi:10.4007/annals.2005.161.489.

\bibitem{KRR2007}
P.~Kurlberg, L.~Rosenzweig and Z.~Rudnick,
\emph{Matrix elements for the quantum cat map: fluctuations in short windows},
arXiv:math/0701685 (2007).

\bibitem{DJ2024}
S.~Dyatlov and M.~J\'ez\'equel,
\emph{Semiclassical Measures for Higher-Dimensional Quantum Cat Maps},
Ann. Henri Poincar\'e \textbf{25} (2024), 1545--1605,
doi:10.1007/s00023-023-01309-x.

\bibitem{KORS2025}
P.~Kurlberg, A.~Ostafe, Z.~Rudnick and I.~E.~Shparlinski,
\emph{On Quantum Ergodicity for Higher Dimensional Cat Maps},
Commun. Math. Phys. \textbf{406} (2025), 174,
doi:10.1007/s00220-025-05350-1.

\bibitem{BS2026}
S.~Bhakta and I.~E.~Shparlinski,
\emph{On quantum ergodicity for higher-dimensional cat maps modulo prime powers},
Bull. London Math. Soc. \textbf{58} (2026),
doi:10.1112/blms.70300.

\bibitem{PW2014}
Z.~Pluha\v r and H.~A.~Weidenm\"uller,
\emph{Universal Quantum Graphs},
Phys. Rev. Lett. \textbf{112} (2014), 144102,
doi:10.1103/PhysRevLett.112.144102.

\bibitem{PW2015}
Z.~Pluha\v r and H.~A.~Weidenm\"uller,
\emph{Quantum graphs and random-matrix theory},
J. Phys. A \textbf{48} (2015), 275102,
doi:10.1088/1751-8113/48/27/275102.

\bibitem{Chowla1970}
S.~Chowla,
\emph{The nonexistence of nontrivial linear relations between the roots of a
certain irreducible equation},
J. Number Theory \textbf{2} (1970), 120--123,
doi:10.1016/0022-314X(70)90012-0.

\bibitem{Wang1982}
K.~Wang,
\emph{On a theorem of S. Chowla},
J. Number Theory \textbf{15} (1982), 1--4,
doi:10.1016/0022-314X(82)90078-6.

\bibitem{Carmichael1913}
R.~D.~Carmichael,
\emph{On the numerical factors of the arithmetic forms $\alpha^n\pm\beta^n$},
Ann. of Math. \textbf{15} (1913--1914), 49--70,
doi:10.2307/1967798.

\bibitem{Yabuta2001}
M.~Yabuta,
\emph{A Simple Proof of Carmichael's Theorem on Primitive Divisors},
Fibonacci Quart. \textbf{39} (2001), 439--443,
doi:10.1080/00150517.2001.12428701.

\bibitem{BHV2001}
Y.~Bilu, G.~Hanrot and P.~M.~Voutier,
\emph{Existence of primitive divisors of Lucas and Lehmer numbers},
J. reine angew. Math. \textbf{539} (2001), 75--122,
doi:10.1515/crll.2001.080.

\bibitem{VM2015}
I.~V.~Vyugin and S.~Makarychev,
\emph{Solutions of polynomial equation over $\mathbf F_p$ and new bounds of additive energy},
arXiv:1504.01354v2 (2015).

\bibitem{MV2019}
S.~Makarychev and I.~Vyugin,
\emph{Solutions of Polynomial Equations in Subgroups of $\mathbf F_p^*$},
Arnold Math. J. \textbf{5} (2019), 105--121,
doi:10.1007/s40598-019-00112-z.

\bibitem{MSS2018}
S.~Macourt, I.~D.~Shkredov and I.~E.~Shparlinski,
\emph{Multiplicative Energy of Shifted Subgroups and Bounds On Exponential Sums with Trinomials in Finite Fields},
Canad. J. Math. \textbf{70} (2018), 1319--1338,
doi:10.4153/CJM-2017-044-2.

\bibitem{RW2022}
M.~Rudnev and J.~Wheeler,
\emph{On incidence bounds with M\"obius hyperbolae in positive characteristic},
Finite Fields Appl. \textbf{78} (2022), 101978,
doi:10.1016/j.ffa.2021.101978.

\bibitem{Michel1995}
P.~Michel,
\emph{Rang moyen de familles de courbes elliptiques et lois de Sato--Tate},
Monatsh. Math. \textbf{120} (1995), 127--136.

\bibitem{SFFK2023}
W.~Sawin, A.~Forey, J.~Fres\'an and E.~Kowalski,
\emph{Quantitative sheaf theory},
J. Amer. Math. Soc. \textbf{36} (2023), 653--726,
doi:10.1090/jams/1008.

\bibitem{Zarhin2013}
Yu.~G.~Zarhin,
\emph{Homomorphisms of abelian varieties over geometric fields of finite characteristic},
J. Inst. Math. Jussieu \textbf{12} (2013), 225--236,
doi:10.1017/S1474748012000655.

\bibitem{GroveSaad2025}
B.~Grove and H.~Saad,
\emph{Hypergeometric distributions and joint families of elliptic curves},
Res. Math. Sci. \textbf{12} (2025), 74,
doi:10.1007/s40687-025-00554-2.

\end{thebibliography}

\end{document}
